% Problems moved from Volume 7, Chapter 13 (Almost every number) in the rewrite to the master plan, decision D-15.
% Their solutions use the ergodicity of the Gauss map and Birkhoff's theorem, which are proved in this volume, Chapter 10.
% NOT YET WIRED into main.tex: the Volume 11 rewrite places them (the problem chapter, ch13-harder.tex, is the natural home).
% Notation, as in Volume 7, Chapter 13: x=[0;a_1,a_2,...] in (0,1), convergents p_n/q_n, cylinder I(a_1,...,a_n) of
% length 1/(q_n(q_n+q_{n-1})) (Volume 7, Problem 13.3(a)), Gauss map G(x)={1/x}, e(t)=e^{2 pi i t}.
% Check before wiring: the labels thm:gaussergodic, thm:birkhoff, thm:cfstats, cor:ergexamples, ch:ergodic of this volume,
% and the Volume 7 numbers quoted (Lemma 5.5 squeeze lemma, Theorem 11.12 Lebesgue on monotone functions,
% Theorem 11.14 fundamental theorem of calculus, Problem 13.16 the Cantor measure).
% Sources to carry over: H. Minkowski, Zur Geometrie der Zahlen, Verhandlungen des III. Internationalen
% Mathematiker-Kongresses, Heidelberg, 1904; A. Denjoy, Sur une fonction réelle de Minkowski, Journal de Mathématiques
% Pures et Appliquées 17 (1938), 105--151; R. Salem, On some singular monotonic functions which are strictly increasing,
% Transactions of the American Mathematical Society 53 (1943), 427--439; J. W. S. Cassels, Colloquium Mathematicum 7
% (1959), 95--101; W. M. Schmidt, Pacific Journal of Mathematics 10 (1960), 661--672.

\begin{exercise}[$\star\star\star$]\label{pb:minkowski}
\emph{Minkowski's question mark.} For irrational $x=[0;a_1,a_2,\dots]\in(0,1)$ let $s_n=a_1+\dots+a_n$ and
\[
?(x)=2\sum_{k\ge1}(-1)^{k+1}2^{-s_k}.
\]
\begin{enumerate}[label=(\alph*)]
\item Show that $?(x)=2\sum_{k=1}^n(-1)^{k+1}2^{-s_k}+(-1)^n2^{-s_n}\,?(G^nx)$, and that $?$ maps the irrationals of the cylinder
$I(a_1,\dots,a_n)$ into an interval of length $2^{-s_n}$.
\item Show that $?$ is strictly increasing on the irrationals of $(0,1)$, and that it extends to a continuous increasing bijection
$[0,1]\to[0,1]$. Compute $?\bigl(\frac{\sqrt5-1}2\bigr)$ and $?(\sqrt2-1)$.
\item Show that for almost every $x$, $\frac{s_n(x)}n\to\infty$.
\item Deduce that $?'(x)=0$ for almost every $x$: the function $?$ is continuous, strictly increasing and singular.
\end{enumerate}
\end{exercise}

\begin{hint}{pb:minkowski}
For (b), compare $x<y$ at the first digit where they differ; the maps $t\mapsto\frac{p_n+tp_{n-1}}{q_n+tq_{n-1}}$ and $t\mapsto(-1)^n2^{-s_n}t$ have the same
monotonicity. For (c), Birkhoff's theorem for the Gauss map applied to $\min(a_1,M)$. For (d), compare $?(I)$ with $\lambda(I)$ for the cylinders around
$x$, and use the squeeze lemma of Volume~7, Lemma~5.5.
\end{hint}

\begin{solution}{pb:minkowski}
\idea $?$ sends the cylinder $I(a_1,\dots,a_n)$, of length about $q_n^{-2}$, onto an interval of length $2^{-s_n}$. For a random number $q_n$ grows exponentially
(Lévy) while $s_n$ grows faster than any multiple of $n$, because the digits have infinite mean. So $?$ squeezes the cylinders around almost every point
superexponentially, and its derivative there is $0$.

\writeup (a) The terms of the series with $k>n$ are $2(-1)^{k+1}2^{-s_k}$, and $s_k=s_n+(a_{n+1}+\dots+a_k)$, where $a_{n+j}=a_j(G^nx)$; factoring out
$(-1)^n2^{-s_n}$ leaves the series of $?(G^nx)$. The series of $?(u)$ alternates with strictly decreasing terms whose first is $2\cdot2^{-a_1(u)}\le1$, so
$0<?(u)<1$ for every irrational $u$. Hence $?$ maps the irrationals of $I(a_1,\dots,a_n)$ into the open interval $J(a_1,\dots,a_n)$ with ends
$\Sigma_n=2\sum_{k\le n}(-1)^{k+1}2^{-s_k}$ and $\Sigma_n+(-1)^n2^{-s_n}$.

(b) For $n=1$ the intervals $J(a)=]2^{-a},2^{1-a}[$, $a\ge1$, are disjoint and in decreasing order as $a$ increases, as are the cylinders $I(a)=]\frac1{a+1},\frac1a[$.
Now let $x<y$ be irrational and $k$ the first index with $a_k(x)\ne a_k(y)$. Both lie in $I=I(a_1,\dots,a_{k-1})$; write $x=\psi(u)$, $y=\psi(v)$ with $u=G^{k-1}x$,
$v=G^{k-1}y$ and $\psi(t)=\frac{p_{k-1}+tp_{k-2}}{q_{k-1}+tq_{k-2}}$, and $?(x)=\Sigma_{k-1}+(-1)^{k-1}2^{-s_{k-1}}?(u)$, likewise for $y$. The map $\psi$ is increasing
if and only if $k-1$ is even (its derivative has the sign of $p_{k-2}q_{k-1}-p_{k-1}q_{k-2}=(-1)^{k-1}$), exactly like $t\mapsto(-1)^{k-1}t$. So it suffices to compare
$u$ and $v$, whose first digits differ: if $a_1(u)>a_1(v)$ then $u<v$ and $?(u)<2^{1-a_1(u)}\le2^{-a_1(v)}<?(v)$. Thus $?$ is strictly increasing on the irrationals.

The intervals $J(w)$, for the words $w$ of length $n$, tile $]0,1[$ up to countably many points (true for $n=1$, and $J(wa)$ is the image of $J(a)$ under the affine map
$t\mapsto\Sigma_n+(-1)^n2^{-s_n}t$, so the $J(wa)$ tile $J(w)$), each contains values of $?$, and each has length $2^{-s_n}\le2^{-n}$. So the image of the irrationals is
dense in $[0,1]$. Extend $?$ by $?(y)=\sup\{?(x):x<y\text{ irrational}\}$ and $?(0)=0$. This is non-decreasing; a jump at some $y$ would leave an open interval containing at
most one value of $?$, contradicting density, so $?$ is continuous, hence a continuous increasing surjection $[0,1]\to[0,1]$, strictly increasing because the
irrationals are dense. Finally $?\bigl(\frac{\sqrt5-1}2\bigr)=2\sum_{k\ge1}(-1)^{k+1}2^{-k}=\frac23$ and $?(\sqrt2-1)=2\sum_{k\ge1}(-1)^{k+1}4^{-k}=\frac25$.

(c) The Gauss measure $\gamma$ is equivalent to $\lambda$ (its density lies between $\frac1{2\log2}$ and $\frac1{\log2}$), $G$ is ergodic for it (Theorem~\ref{thm:gaussergodic}),
and $a_{k+1}=a_1\circ G^k$. For $M\ge1$, Birkhoff's theorem (Theorem~\ref{thm:birkhoff}) applied to the bounded function $\min(a_1,M)$ gives, for almost every $x$,
\[
\liminf_n\frac{s_n}n\ge\lim_n\frac1n\sum_{k<n}\min(a_1,M)(G^kx)=\int\min(a_1,M)\dd\gamma.
\]
As $M\to\infty$ the right side tends to $\sum_kk\,\gamma(a_1=k)=\sum_kk\log_2\bigl(1+\frac1{k(k+2)}\bigr)=\infty$, by monotone convergence, since the terms are
equivalent to $\frac1{k\log2}$.

(d) Let $x$ be irrational, satisfying (c) and $q_n(x)\le e^{2n}$ for large $n$ (Theorem~\ref{thm:cfstats}(c)): almost every $x$ does. Let $[u_n,v_n]$ be the closure of the
cylinder of rank $n$ containing $x$; then $u_n<x<v_n$, $v_n-u_n=\frac1{q_n(q_n+q_{n-1})}\ge\frac1{2q_n^2}$, and $?(v_n)-?(u_n)\le2^{-s_n}$ by (a) and continuity. So
\[
\frac{?(v_n)-?(u_n)}{v_n-u_n}\le2q_n^22^{-s_n}\le2\exp\bigl(4n-s_n\log2\bigr)\longrightarrow0.
\]
By Lebesgue's theorem on monotone functions (Volume~7, Theorem~11.12), the monotone function $?$ is differentiable almost everywhere, and at a point of differentiability with the properties above
the squeeze lemma (Volume~7, Lemma~5.5) gives $?'(x)=0$. If $?$ were absolutely continuous we would have $1=?(1)-?(0)=\int_0^1?'=0$ (Volume~7, Theorem~11.14).

\tools{Volume~2, Chapter~9; Volume~7, Chapters~5 and~11 (Lebesgue's theorem on monotone functions); Chapter~\ref{ch:ergodic} (Birkhoff's theorem for the Gauss map).}
\leads Minkowski introduced $?$ in 1904 because it sends the quadratic irrationals to the rationals, as (b) shows on two examples, and the rationals to the dyadic
rationals. Its singularity was proved by Denjoy (1938) and Salem (1943). There is a probabilistic reading: $?$ is the distribution function of $[0;A_1,A_2,\dots]$ with
independent digits $P(A_i=k)=2^{-k}$, since the cylinder $I(a_1,\dots,a_n)$ has $?$-measure $2^{-s_n}$. Under this law the digits have mean $2$, under Lebesgue measure
mean $\infty$, and the strong law of large numbers separates the two measures. Figure~\ref{fig:minkowski} shows the graph.
\end{solution}

\begin{figure}[ht]
\centering
\fig{from-volume7-minkowski}
\caption{Minkowski's question mark $?(x)$ (charcoal) against the diagonal (grey). It is continuous and strictly increasing, fixes $0$, $\frac12$ and $1$, and sends
$\frac{\sqrt5-1}2=0.618\ldots$ to $\frac23$ and $\sqrt2-1=0.414\ldots$ to $\frac25$ (rust dots); yet its derivative vanishes almost everywhere
(Problem~\ref{pb:minkowski}).}\label{fig:minkowski}
\end{figure}

\begin{exercise}[$\star\star$]\label{pb:cantororbit}
\emph{The orbit of a Cantor point under $\times3$.} Let $\mu$ be the Cantor measure, the law of $X=\sum_{k\ge1}2\varepsilon_k3^{-k}$ with $\varepsilon_k$ independent and
$P(\varepsilon_k=0)=P(\varepsilon_k=1)=\frac12$ (Volume~7, Problem~13.16, where it is shown that $(3^nx)$ is not equidistributed modulo $1$ for any $x$ in the Cantor set).
Show that for $\mu$-almost every $x$ the sequence $(3^nx)$ is distributed according to $\mu$: for every
continuous $g$ on $[0,1]$ with $g(0)=g(1)$, $\frac1N\sum_{n<N}g(\{3^nx\})\to\int g\dd\mu$.
\end{exercise}

\begin{hint}{pb:cantororbit}
The map $x\mapsto3x\bmod1$ shifts the ternary digits; it preserves $\mu$, and Birkhoff's theorem applies.
\end{hint}

\begin{solution}{pb:cantororbit}
\idea Multiplication by $3$ is the shift on the ternary digits, and the Cantor measure is the Bernoulli measure on them.

\writeup Let $T(x)=\{3x\}$. On ternary digits $T$ is the shift, and $\mu$ is the image of the product measure on $\{0,2\}^{\N^*}$, so $\mu$ is $T$-invariant and $T$ is
ergodic for $\mu$ (Corollary~\ref{cor:ergexamples}(a)). Let $(g_j)$ be a countable dense subset of the continuous $g$ with $g(0)=g(1)$, for the sup norm. By Birkhoff's theorem
(Theorem~\ref{thm:birkhoff}), for $\mu$-almost every $x$, $\frac1N\sum_{n<N}g_j(T^nx)\to\int g_j\dd\mu$ for all $j$, hence for all such $g$ by uniform approximation; and $T^nx=\{3^nx\}$.

\tools{Chapter~\ref{ch:ergodic} (Birkhoff's theorem, Bernoulli shifts); Volume~7, Chapter~13.}
\leads Cassels (1959) and Schmidt (1960) proved that $\mu$-almost every $x$ is nevertheless normal in base $2$: the Cantor measure is invisible to $\times3$ but
looks like Lebesgue measure to $\times2$.
\end{solution}
